\myGraphic{0.227}{images/Mak_Author-Photo.png}{}

Ronald Mak（罗纳德·马克）是硅谷圣何塞州立大学备受好评的讲师，教授 C++、Java 和 Python 的面向对象分析与设计。作为 NASA 和 JPL 的资深计算机科学家，他为多项重大太空任务开发软件，包括火星漫游车、国际空间站、立方星和猎户座飞船。他曾是 IBM 研究院的研究人员，在那里协助为一个超级计算机开发系统软件，并为一个先进的数据科学项目开发软件。他曾担任劳伦斯利弗莫尔国家实验室（Lawrence Livermore National Laboratory）国家点火装置（National Ignition Facility，NIF）的企业软件策略师，该装置正在聚变能源领域取得重要突破。在职业生涯早期，他曾在多家硅谷公司担任软件开发者和工程经理，包括惠普（Hewlett-Packard）、太阳微系统（Sun Microsystems）、苹果（Apple）以及几家初创公司。他拥有斯坦福大学数学科学和计算机科学学位。他是七项软件专利的发明人。他撰写过关于编译器编写、软件工程和数值计算的书，这些书已被翻译成多种语言。尽管 Ron 曾研究行星的相对运动，并做过涉及爱因斯坦相对论的计算，他至今仍然惊叹于太阳每天早晨照常升起、自行车不会翻倒。

关于技术编辑

Juan Rufes 写代码至今已有 30 多年。他从一台 8 位 MSX 家用电脑起步，后来转向杀毒和加密方面的工作，最近 12 年则在金融领域，主要用 C++ 开发低延迟交易系统，偶尔也用一些 C。Juan 喜欢底层和系统级编程，尤其喜欢 C++ 代码优化。在业余时间，他喜欢读一本好书，再喝上一两杯好茶。
